% Generated by scripts/assemble.py; corrections belong in source/corrections.json.
\examhead{25}{中国矿业大学(北京) 2024 年数学分析试卷}{29}
% p030-b06; source PDF page 30
\blockanchor{p030-b06}
\question{1}(10 分) 计算函数极限 $\lim_{x\to 0}\frac{(1 + x)^{\frac{2}{x}} - e^2(1 - \ln(1 + x))}{x^2}.$\par
% p030-b07; source PDF page 30
\blockanchor{p030-b07}
\question{2}(10 分) 计算不定积分 $\int \frac{xe^x}{\sqrt{e^x + 1}} \, \mathrm{d}x$.\par
% p030-b08; source PDF page 30
\blockanchor{p030-b08}
\question{3}(10 分) 函数 $z = z(x, y)$ 是由方程 $\frac{x}{z} = \ln \frac{z}{y}$ 确定的隐函数, 计算 $\frac{\partial^2 z}{\partial x \partial y}$.\par
% p030-b09; source PDF page 30
\blockanchor{p030-b09}
\question{4}(10 分) 计算二重积分 $\iint_{D} e^{-(x + y)^2} \, \mathrm{d}x \, \mathrm{d}y$, 其中 $D$ 是由 $x + y = 1, x = 0, y = x$ 所围成的区域.\par
% p030-b10; source PDF page 30
\blockanchor{p030-b10}
\question{5}(10 分) 计算曲线积分 $\oint_{L} \frac{3x \, \mathrm{d}x - 4y \, \mathrm{d}y}{3x^{2} + 4y^{2}}$, 其中 $L$ 是 $x^{2} + y^{2} = 1$ 的正向边界.\par
% p030-b11; source PDF page 30
\par\Needspace{6\baselineskip}
\blockanchor{p030-b11}
\question{6}(10 分) 计算曲面积分\par
\nopagebreak[4]
% p030-b12; source PDF page 30
\blockanchor{p030-b12}
\displayformula{\iint_ {\Sigma} y ^ {3} \mathrm{d} z \mathrm{d} x + (y + z) \mathrm{d} x \mathrm{d} y.}
% p030-b13; source PDF page 30
\blockanchor{p030-b13}
其中 $\Sigma$ 为 $z = x^{2} + y^{2}$ 被 $z = 0, z = 1$ 所截部分的外侧.\par
% p030-b14; source PDF page 30
\blockanchor{p030-b14}
\question{7}(15 分) 求函数 $f(x) = \arctan \frac{1 - x^2}{1 + x^2}$ 的单调区间, 极值, 凹凸区间, 拐点和渐近线.\par
% p030-b15; source PDF page 30
\blockanchor{p030-b15}
\question{8}(15 分) 函数 $f(x, y) = x^2 + (y - 1)^2 (x \neq 0)$ 在 $\frac{x^2}{a^2} + \frac{y^2}{b^2} = 1$ 条件下有最小值 1，且椭圆的面积最小，求 $a, b$ 的值。(可能有误)\par
% p030-b16; source PDF page 30
\blockanchor{p030-b16}
\question{9}(15 分) 判断反常积分 $\int_{0}^{+\infty} \frac{e^{\sin x} \sin 2x}{x^p} \, \mathrm{d}x$ ($p > 0$) 的敛散性, 并讨论条件收敛和绝对收敛.\par
% p030-b17; source PDF page 30
\blockanchor{p030-b17}
\question{10}(15 分) 函数 $f(x), g(x)$ 在区间 $[a, b]$ 上二阶可导, 且 $f(a) = g(a), f(b) = g(b), f(x), g(x)$ 在 $(a, b)$ 内有相同的最大值. 证明: 存在 $\xi \in (a, b)$, 使得 $f''(\xi) = g''(\xi)$.\par
% p030-b18; source PDF page 30
\blockanchor{p030-b18}
\question{11}(15 分) 设 $f(x) = \sum_{n=1}^{\infty} \frac{\cos nx}{\sqrt{n^3 + n}}$, $g(x)$ 是 $f(x)$ 的原函数, 且 $g(0) = 0$, 证明:\par
% p030-b19; source PDF page 30
\blockanchor{p030-b19}
(1) $f(x)$ 在 $\mathbb{R}$ 上连续.\par
% p030-b20; source PDF page 30
\blockanchor{p030-b20}
(2) $\frac{\sqrt{2}}{2} - \frac{1}{15} < g\left(\frac{\pi}{2}\right) < \frac{\sqrt{2}}{2}$.\par
% p030-b21; source PDF page 30
\blockanchor{p030-b21}
\question{12}(15 分) 叙述实数系的闭区间套定理和确界原理, 并用闭区间套定理证明确界原理.\par
